%% Déphasage entre tension et intensité en régime sinusoïdal.
%% u(t) : cosinusoïde d'amplitude 1, période 1 s, maximale en t = 0, 1 et 2 s.
%% i(t) : amplitude 0,5, même période, retardée de 0,1 s (phi = 36°).
\documentclass[tikz,border=4pt]{standalone}
\usepackage{liaisons}
\begin{document}
\begin{tikzpicture}[x=2.6cm, y=1.15cm,
  axe/.style={-{Stealth[length=5pt]}, line width=0.6pt},
  grille/.style={black!12, line width=0.3pt},
  grad/.style={font=\scriptsize},
  cote/.style={{Stealth[length=4pt]}-{Stealth[length=4pt]}, line width=0.6pt}]

%% ---- quadrillage ---------------------------------------------------------
\foreach \gx in {0,0.125,...,2.51} { \draw[grille] (\gx,-1.08) -- (\gx,1.08); }
\foreach \gy in {-1,-0.875,...,1.01} { \draw[grille] (0,\gy) -- (2.5,\gy); }

%% ---- axes et graduations -------------------------------------------------
\draw[axe] (-0.06,0) -- (2.72,0) node[anchor=south east, inner sep=3pt, font=\small] {$t$ (s)};
\draw[axe] (0,-1.3) -- (0,1.35);
\foreach \t/\lab in {0.5/{0,5}, 1/{1}, 1.5/{1,5}, 2/{2}, 2.5/{2,5}} {
  \draw[line width=0.5pt] (\t,0) -- (\t,-0.06) node[grad, anchor=north, inner sep=2pt] {\lab}; }
\node[grad, anchor=north east, inner sep=2pt] at (0,0) {0};
\foreach \v/\lab in {-1/{$-1$}, -0.5/{$-0{,}5$}, 0.5/{$0{,}5$}, 1/{$1$}} {
  \draw[line width=0.5pt] (0,\v) -- (-0.015,\v) node[grad, anchor=east, inner sep=2pt] {\lab}; }

%% ---- les deux courbes ----------------------------------------------------
\draw[solideB, line width=1.3pt] plot[domain=0:2.5, samples=180] (\x,{cos(\x*360)});
\draw[solideA, line width=1.3pt] plot[domain=0:2.5, samples=180] (\x,{0.5*cos((\x-0.1)*360)});
\node[font=\footnotesize, text=solideB, anchor=south, inner sep=2pt] at (2,1.02) {$u(t)$};
\node[font=\footnotesize, text=solideA, anchor=west, inner sep=2pt] at (2.16,0.52) {$i(t)$};

%% ---- déphasage entre les deux maxima ------------------------------------
\draw[black!45, line width=0.4pt, dotted] (1,1) -- (1,1.28);
\draw[black!45, line width=0.4pt, dotted] (1.1,0.5) -- (1.1,1.28);
\draw[cote] (1,1.22) -- (1.1,1.22);
\node[font=\footnotesize, anchor=west, inner sep=3pt] at (1.13,1.22) {$\varphi$};

%% ---- expressions des deux signaux ---------------------------------------
\node[anchor=north west, font=\footnotesize, text=solideB, inner sep=2pt] at (0,-1.42)
  {$u(t) = U\sqrt{2}\cos(\omega t)$};
\node[anchor=north west, font=\footnotesize, text=solideA, inner sep=2pt] at (0,-1.75)
  {$i(t) = I\sqrt{2}\cos(\omega t - \varphi)$};
\end{tikzpicture}
\end{document}
